\documentclass[a4paper,14pt]{extarticle}
\usepackage[UTF8]{ctex}
\usepackage[T2A]{fontenc}
\usepackage[russian]{babel}
\usepackage{geometry}
\usepackage{graphicx}
\usepackage{amsmath}
\usepackage{amsfonts}
\usepackage{amssymb}
\usepackage{hyperref}
\usepackage{caption}
\usepackage{subcaption}
\usepackage{listings}
\usepackage{multirow}
\usepackage{xcolor}
\usepackage{float}

\geometry{left=30mm, right=15mm, top=20mm, bottom=20mm}

\definecolor{codegreen}{rgb}{0,0.6,0}
\definecolor{codegray}{rgb}{0.5,0.5,0.5}
\definecolor{codepurple}{rgb}{0.58,0,0.82}
\definecolor{backcolour}{rgb}{0.95,0.95,0.92}

\lstset{
    inputencoding=utf8,
    extendedchars=true,
    literate={а}{{\cyra}}1 {б}{{\cyrb}}1 {в}{{\cyrv}}1 {г}{{\cyrg}}1 {д}{{\cyrd}}1 {е}{{\cyre}}1 {ё}{{\cyryo}}1 {ж}{{\cyrzh}}1 {з}{{\cyrz}}1 {и}{{\cyri}}1 {й}{{\cyrishrt}}1 {к}{{\cyrk}}1 {л}{{\cyrl}}1 {м}{{\cyrm}}1 {н}{{\cyrn}}1 {о}{{\cyro}}1 {п}{{\cyrp}}1 {р}{{\cyrr}}1 {с}{{\cyrs}}1 {т}{{\cyrt}}1 {у}{{\cyru}}1 {ф}{{\cyrf}}1 {х}{{\cyrkh}}1 {ц}{{\cyrc}}1 {ч}{{\cyrch}}1 {ш}{{\cyrsh}}1 {щ}{{\cyrshch}}1 {ъ}{{\cyrhrdsn}}1 {ы}{{\cyrery}}1 {ь}{{\cyrsftsn}}1 {э}{{\cyre}}1 {ю}{{\cyryu}}1 {я}{{\cyrya}}1 {А}{{\cyra}}1 {Б}{{\cyrb}}1 {В}{{\cyrv}}1 {Г}{{\cyrg}}1 {Д}{{\cyrd}}1 {Е}{{\cyre}}1 {Ё}{{\cyryo}}1 {Ж}{{\cyrzh}}1 {З}{{\cyrz}}1 {И}{{\cyri}}1 {Й}{{\cyrishrt}}1 {К}{{\cyrk}}1 {Л}{{\cyrl}}1 {М}{{\cyrm}}1 {Н}{{\cyrn}}1 {О}{{\cyro}}1 {П}{{\cyrp}}1 {Р}{{\cyrr}}1 {С}{{\cyrs}}1 {Т}{{\cyrt}}1 {У}{{\cyru}}1 {Ф}{{\cyrf}}1 {Х}{{\cyrkh}}1 {Ц}{{\cyrc}}1 {Ч}{{\cyrch}}1 {Ш}{{\cyrsh}}1 {Щ}{{\cyrshch}}1 {Ъ}{{\cyrhrdsn}}1 {Ы}{{\cyrery}}1 {Ь}{{\cyrsftsn}}1 {Э}{{\cyre}}1 {Ю}{{\cyryu}}1 {Я}{{\cyrya}}1,
    basicstyle=\ttfamily\footnotesize,
    breaklines=true,
    numbers=left,
    numberstyle=\tiny,
    stepnumber=1,
    numbersep=5pt,
    backgroundcolor=\color{gray!5},
    showspaces=false,
    showstringspaces=false,
    showtabs=false,
    frame=single,
    tabsize=2,
    captionpos=b,
    breakatwhitespace=false,
    language=Python
}

\begin{document}

\begin{titlepage}
\begin{center}
\large
Санкт-Петербургский политехнический университет \\
Петра Великого \\
Высшая школа прикладной математики и вычислительной физики \\
\vspace{3cm}

\textbf{Отчёт по лабораторной работе по дисциплине} \\
«Математическая статистика» \\
\vspace{0.5cm}

\textbf{ЛАБОРАТОРНАЯ РАБОТА №4} \\
«Эмпирическая функция распределения и ядерные оценки плотности» \\
\vspace{2cm}

Выполнил: Гуцалов Егор Никитич \\
Группа: 5030102/30201 \\
\vfill

Санкт-Петербург \\
2025
\end{center}
\end{titlepage}

\section*{2. Постановка задачи}

\textbf{Цель работы:} Изучить методы построения эмпирической функции распределения и ядерных оценок плотности, проанализировать влияние размера выборки на точность оценок, сравнить ядерные оценки с гистограммами.

\textbf{Задачи:}
\begin{itemize}
    \item Сгенерировать выборки объёмом $n = 20, 60, 100$ для каждого из 5 распределений
    \item Построить эмпирическую функцию распределения (ЭФР) и теоретическую функцию распределения
    \item Построить ядерную оценку плотности (с гауссовым ядром) и теоретическую плотность
    \item Выполнить построение на отрезке $[-4, 4]$ для непрерывных распределений и на отрезке $[6, 14]$ для распределения Пуассона
    \item Проанализировать, как мощность выборки влияет на приближение эмпирической функции к истинной
    \item Сравнить два способа оценки плотности: ядерные оценки и гистограммы
\end{itemize}

\textbf{Исследуемые распределения:}
\begin{enumerate}
    \item Нормальное $N(0, 1)$
    \item Коши $C(0, 1)$
    \item Лапласа $L(0, \frac{1}{\sqrt{2}})$
    \item Пуассона $P(10)$
    \item Равномерное $U(-\sqrt{3}, \sqrt{3})$
\end{enumerate}

\section*{3. Теоретическая часть}

\subsection*{3.1 Эмпирическая функция распределения}

\textbf{Эмпирическая функция распределения} (ЭФР) — это ступенчатая функция, определяемая как доля элементов выборки, не превосходящих заданного значения $x$:

\begin{equation}
\hat{F}_n(x) = \frac{1}{n} \sum_{i=1}^{n} I(x_i \leq x)
\end{equation}

где $I(\cdot)$ — индикаторная функция, равная 1, если условие выполняется, и 0 иначе.

\textbf{Свойства ЭФР:}
\begin{itemize}
    \item Неубывающая ступенчатая функция
    \item Принимает значения от 0 до 1
    \item Делает скачок величины $1/n$ в каждой точке выборки
    \item Состоятельная оценка истинной функции распределения
\end{itemize}

\textbf{Теорема Гливенко-Кантелли:}
\begin{equation}
\sup_{x} |\hat{F}_n(x) - F(x)| \xrightarrow[n \to \infty]{a.s.} 0
\end{equation}

Эта теорема утверждает, что эмпирическая функция распределения равномерно сходится к истинной функции распределения почти наверное.

\subsection*{3.2 Ядерная оценка плотности}

\textbf{Ядерная оценка плотности} (Kernel Density Estimation, KDE) — это непараметрический метод оценки плотности вероятности, основанный на сглаживании данных.

В отличие от гистограммы, которая является ступенчатой функцией, ядерная оценка даёт \textbf{непрерывную гладкую оценку} плотности.

\textbf{Формула ядерной оценки:}
\begin{equation}
\hat{f}_h(x) = \frac{1}{n} \sum_{i=1}^{n} K_h(x - x_i) = \frac{1}{nh} \sum_{i=1}^{n} K\left(\frac{x - x_i}{h}\right)
\end{equation}

где:
\begin{itemize}
    \item $K(\cdot)$ — ядерная функция (в данной работе используется гауссово ядро)
    \item $h > 0$ — параметр сглаживания (bandwidth)
\end{itemize}

\textbf{Гауссово ядро:}
\begin{equation}
K(u) = \frac{1}{\sqrt{2\pi}} e^{-u^2/2}
\end{equation}

\textbf{Выбор параметра сглаживания $h$:}
\begin{itemize}
    \item Малое $h$ — оценка следует за данными, может быть шумной (переобучение)
    \item Большое $h$ — оценка слишком гладкая, может потерять детали (недообучение)
    \item В работе используется автоматический выбор $h$ по правилу Скотта
\end{itemize}

\subsection*{3.3 Сравнение гистограммы и KDE}

\begin{table}[H]
\centering
\caption{Сравнение гистограммы и ядерной оценки плотности}
\begin{tabular}{|l|p{5cm}|p{5cm}|}
\hline
\textbf{Характеристика} & \textbf{Гистограмма} & \textbf{KDE} \\
\hline
Гладкость & Ступенчатая функция & Непрерывная гладкая \\
\hline
Зависимость от параметров & Выбор границ и ширины бинов & Выбор параметра сглаживания $h$ \\
\hline
Визуальное качество & Грубая, дискретная & Плавная, естественная \\
\hline
Вычислительная сложность & Низкая & Выше \\
\hline
Интерпретация & Простая & Требует понимания ядер \\
\hline
\end{tabular}
\end{table}

\section*{4. Реализация}

\textbf{Язык программирования:} Python 3.x

\textbf{Используемые библиотеки:}
\begin{itemize}
    \item \texttt{numpy} — генерация случайных выборок и вычисления
    \item \texttt{matplotlib.pyplot} и \texttt{seaborn} — визуализация данных
    \item \texttt{scipy.stats} — теоретические распределения и гауссово ядро
\end{itemize}

\textbf{Среда разработки:} Visual Studio Code

\subsection*{Основные этапы реализации:}

\begin{enumerate}
    \item Генерация выборок объёмом $n = 20, 60, 100$ из каждого распределения
    \item Построение эмпирической функции распределения (ступенчатый график)
    \item Наложение теоретической функции распределения
    \item Построение гистограммы выборки
    \item Вычисление ядерной оценки плотности с гауссовым ядром
    \item Наложение теоретической плотности
    \item Сохранение графиков для всех распределений и размеров выборок
\end{enumerate}

\begin{lstlisting}[caption={Построение ЭФР и ядерной оценки плотности}]
import numpy as np
import matplotlib.pyplot as plt
from scipy import stats
from scipy.stats import gaussian_kde

n = 100
sample = np.random.normal(0, 1, n)

x_ecdf = np.linspace(-4, 4, 1000)
y_ecdf = np.array([np.mean(sample <= xi) 
                   for xi in x_ecdf])

kde = gaussian_kde(sample)
y_kde = kde(x_ecdf)

y_theoretical_pdf = stats.norm.pdf(
    x_ecdf, loc=0, scale=1)

plt.figure(figsize=(10, 6))
plt.hist(sample, bins='sqrt', density=True, 
         alpha=0.3, label='Histogram')
plt.plot(x_ecdf, y_kde, 'b-', linewidth=2, 
         label='Kernel Density')
plt.plot(x_ecdf, y_theoretical_pdf, 'r--', 
         linewidth=2, label='Theoretical')
plt.legend()
plt.show()
\end{lstlisting}

\section*{5. Результаты}

\subsection*{5.1 Нормальное распределение $N(0, 1)$}

\begin{figure}[H]
\centering
\includegraphics[width=0.9\textwidth]{results_lab4/normal_ecdf.png}
\caption{Эмпирическая функция распределения для нормального закона}
\label{fig:normal_ecdf}
\end{figure}

\begin{figure}[H]
\centering
\includegraphics[width=0.9\textwidth]{results_lab4/normal_kde.png}
\caption{Ядерная оценка плотности для нормального закона}
\label{fig:normal_kde}
\end{figure}

\textbf{Наблюдения:}
\begin{itemize}
    \item При $n=20$ ЭФР заметно отклоняется от теоретической кривой
    \item При $n=100$ наблюдается хорошее совпадение ЭФР с теорией
    \item KDE плавно приближается к теоретической плотности при росте $n$
    \item Гистограмма при малых $n$ имеет разреженные бины
\end{itemize}

\subsection*{5.2 Распределение Коши $C(0, 1)$}

\begin{figure}[H]
\centering
\includegraphics[width=0.9\textwidth]{results_lab4/cauchy_ecdf.png}
\caption{Эмпирическая функция распределения для закона Коши}
\label{fig:cauchy_ecdf}
\end{figure}

\begin{figure}[H]
\centering
\includegraphics[width=0.9\textwidth]{results_lab4/cauchy_kde.png}
\caption{Ядерная оценка плотности для закона Коши}
\label{fig:cauchy_kde}
\end{figure}

\textbf{Наблюдения:}
\begin{itemize}
    \item Из-за тяжёлых хвостов наблюдаются значительные флуктуации
    \item KDE сглаживает выбросы, но может терять детали распределения
    \item Сходимость медленнее, чем для нормального распределения
\end{itemize}

\subsection*{5.3 Распределение Лапласа $L(0, 1/\sqrt{2})$}

\begin{figure}[H]
\centering
\includegraphics[width=0.9\textwidth]{results_lab4/laplace_ecdf.png}
\caption{Эмпирическая функция распределения для закона Лапласа}
\label{fig:laplace_ecdf}
\end{figure}

\begin{figure}[H]
\centering
\includegraphics[width=0.9\textwidth]{results_lab4/laplace_kde.png}
\caption{Ядерная оценка плотности для закона Лапласа}
\label{fig:laplace_kde}
\end{figure}

\textbf{Наблюдения:}
\begin{itemize}
    \item Острый пик в центре хорошо виден на KDE
    \item ЭФР демонстрирует более быстрый рост в центре
    \item При $n=100$ оценка плотности близка к теоретической
\end{itemize}

\subsection*{5.4 Распределение Пуассона $P(10)$}

\begin{figure}[H]
\centering
\includegraphics[width=0.9\textwidth]{results_lab4/poisson_ecdf.png}
\caption{Эмпирическая функция распределения для закона Пуассона}
\label{fig:poisson_ecdf}
\end{figure}

\begin{figure}[H]
\centering
\includegraphics[width=0.9\textwidth]{results_lab4/poisson_kde.png}
\caption{Ядерная оценка плотности для закона Пуассона}
\label{fig:poisson_kde}
\end{figure}

\textbf{Наблюдения:}
\begin{itemize}
    \item Дискретная природа распределения видна на ступенчатой ЭФР
    \item KDE применяет сглаживание к дискретным данным
    \item При $\lambda=10$ распределение приближается к нормальному
\end{itemize}

\subsection*{5.5 Равномерное распределение $U(-\sqrt{3}, \sqrt{3})$}

\begin{figure}[H]
\centering
\includegraphics[width=0.9\textwidth]{results_lab4/uniform_ecdf.png}
\caption{Эмпирическая функция распределения для равномерного закона}
\label{fig:uniform_ecdf}
\end{figure}

\begin{figure}[H]
\centering
\includegraphics[width=0.9\textwidth]{results_lab4/uniform_kde.png}
\caption{Ядерная оценка плотности для равномерного закона}
\label{fig:uniform_kde}
\end{figure}

\textbf{Наблюдения:}
\begin{itemize}
    \item Линейный рост ЭФР характерен для равномерного распределения
    \item KDE размывает резкие границы из-за сглаживания
    \item При больших $n$ оценка приближается к прямоугольной форме
\end{itemize}

\subsection*{5.6 Сравнение всех распределений}

\begin{figure}[H]
\centering
\includegraphics[width=0.9\textwidth]{results_lab4/all_distributions_comparison.png}
\caption{Сравнение ЭФР и KDE для всех распределений ($n=100$)}
\label{fig:all_comparison}
\end{figure}

\section*{6. Обсуждение}

\subsection*{6.1 Влияние размера выборки}

\begin{enumerate}
    \item \textbf{Эмпирическая функция распределения:}
    \begin{itemize}
        \item При $n=20$: заметные случайные отклонения от теоретической кривой
        \item При $n=60$: улучшение сходимости, уменьшение флуктуаций
        \item При $n=100$: хорошее согласие с теорией, подтверждение теоремы Гливенко-Кантелли
    \end{itemize}
    
    \item \textbf{Ядерная оценка плотности:}
    \begin{itemize}
        \item При малых $n$: оценка может быть шумной или слишком сглаженной
        \item При увеличении $n$: KDE всё лучше воспроизводит форму теоретической плотности
        \item Автоматический выбор параметра сглаживания работает удовлетворительно
    \end{itemize}
\end{enumerate}

\subsection*{6.2 Сравнение методов оценки плотности}

\textbf{Преимущества KDE перед гистограммой:}
\begin{itemize}
    \item Непрерывность и гладкость оценки
    \item Отсутствие зависимости от выбора границ бинов
    \item Лучшее визуальное восприятие формы распределения
    \item Более эффективное использование данных
\end{itemize}

\textbf{Недостатки KDE:}
\begin{itemize}
    \item Вычислительная сложность выше
    \item Требует выбора параметра сглаживания
    \item Может размывать резкие границы (для равномерного распределения)
    \item Менее интуитивна для интерпретации
\end{itemize}

\subsection*{6.3 Особенности различных распределений}

\begin{enumerate}
    \item \textbf{Нормальное:} быстрая сходимость, обе оценки работают хорошо
    \item \textbf{Коши:} медленная сходимость из-за тяжёлых хвостов, KDE сглаживает выбросы
    \item \textbf{Лапласа:} острый пик хорошо воспроизводится при достаточном $n$
    \item \textbf{Пуассона:} дискретная природа требует осторожности при интерпретации KDE
    \item \textbf{Равномерное:} KDE размывает границы, гистограмма лучше показывает прямоугольную форму
\end{enumerate}

\section*{7. Выводы}

\begin{enumerate}
    \item Освоены методы построения эмпирической функции распределения и ядерных оценок плотности для анализа одномерных данных.
    
    \item Подтверждена теорема Гливенко-Кантелли: с ростом объёма выборки эмпирическая функция распределения равномерно сходится к истинной функции распределения.
    
    \item Показано, что ядерная оценка плотности даёт более гладкую и визуально приятную оценку по сравнению с гистограммой, не зависящую от выбора границ бинов.
    
    \item Установлено, что скорость сходимости оценок зависит от типа распределения:
    \begin{itemize}
        \item Быстрая сходимость для нормального и равномерного распределений
        \item Медленная сходимость для распределения Коши из-за тяжёлых хвостов
    \end{itemize}
    
    \item Выявлено, что для дискретных распределений (Пуассона) применение KDE требует осторожности в интерпретации, так как метод предполагает непрерывность.
    
    \item Приобретены практические навыки сравнительного анализа различных методов оценки плотности вероятности в Python.
\end{enumerate}

\section*{8. Список литературы}

\begin{enumerate}
    \item Гмурман В.Е. Теория вероятностей и математическая статистика. — М.: Высшая школа, 2003.
    
    \item Кобецкая И.А. Математическая статистика: Учебное пособие. — СПб.: Лань, 2019.
    
    \item Silverman B.W. Density Estimation for Statistics and Data Analysis. — Chapman and Hall, 1986.
    
    \item SciPy documentation. URL: \url{https://docs.scipy.org/doc/scipy/}
    
    \item Matplotlib documentation. URL: \url{https://matplotlib.org/}
    
    \item Seaborn documentation. URL: \url{https://seaborn.pydata.org/}
\end{enumerate}

\section*{9. Приложение}

Исходный код программы размещён в репозитории GitHub:

\url{https://github.com/Ftp-Glich/Statistics}

В папке \texttt{results\_lab4} находятся:
\begin{itemize}
    \item \texttt{[distribution]\_ecdf.png} — эмпирические функции распределения (5 файлов)
    \item \texttt{[distribution]\_kde.png} — ядерные оценки плотности (5 файлов)
    \item \texttt{all\_distributions\_comparison.png} — сравнительный график для всех распределений
\end{itemize}

\end{document}